% TikZ 图：光电效应伏安特性
% Day 27：近代物理
\begin{tikzpicture}[scale=0.9,>=Latex]
  \draw[->,thick] (-2.4,0) -- (4.0,0) node[right] {$U$};
  \draw[->,thick] (0,0) -- (0,3.0) node[above] {$I$};
  \coordinate (Uc) at (-1.35,0);
  \draw[blue,thick] (Uc) .. controls (-0.72,1.08) and (0.45,2.15) .. (3.3,2.25);
  \draw[red,thick] (Uc) .. controls (-0.55,0.72) and (0.65,1.48) .. (3.3,1.62);
  \draw[dashed] (Uc) -- (-1.35,-0.34);
  \node[below] at (-1.35,-0.34) {$-U_c$};
  \node[blue] at (3.55,2.25) {强光};
  \node[red] at (3.55,1.65) {弱光};
  \node[font=\small,align=center] at (2.05,2.75) {同频率、不同光强};
  \node[font=\small,align=center] at (1.65,-1.05) {同频率时截止电压相同；饱和电流反映光电子数};
\end{tikzpicture}
